\section{Задание}

В рамках выполнения лабораторной работы \textit{по варианту №1} необходимо выполнить следующие операции:
\begin{enumerate}
  \item Сконфигурировать в своём домашнем каталоге репозитории \texttt{Git} и \texttt{Subversion}.
  \item Воспроизвести последовательность команд для систем контроля версий \texttt{Git} и \texttt{Subversion}, которые осуществляют операции над исходным кодом в соответствии с блок-схемой:
  \begin{center}
    \begin{adjustbox}{max width=\textwidth}
      \input{tikz/Plot}
    \end{adjustbox}
  \end{center}
\end{enumerate}

При составлении последовательности команд необходимо учитывать следующие условия:
\begin{itemize}
  \item Цвет элементов схемы указывает на пользователя, который совершил действие \textit{(красный --- первый, синий --- второй)}.
  \item Цифры над узлами --- номер ревизии. Ревизии создаются последовательно.
  \item Необходимо разрешать конфликты между версиями, если при слиянии они возникают.
\end{itemize}